\section{Digital Twin Simulation}
\label{sec:digital-twin}

The digital twin provides a physics-informed simulation layer that validates
copilot recommendations before they are presented to the engineering team.

\subsection{Model Structure}

The digital twin is composed of three sub-models:

\begin{description}
  \item[Tyre degradation model] predicts surface and carcass temperature
    evolution as a function of load, speed, and ambient conditions.
  \item[Suspension kinematics model] maps setup changes (ride height, spring
    rate, damper settings) to predicted handling characteristics.
  \item[Lap time delta model] aggregates tyre and suspension effects into a
    predicted lap time change for each proposed setup modification.
\end{description}

\subsection{Validation Protocol}

Before a recommendation is released to the \rcc, the digital twin executes a
validation run:

\begin{enumerate}
  \item Apply the proposed setup change to the current bike state.
  \item Simulate the next five laps under projected conditions.
  \item If predicted lap time delta is negative (improvement): pass.
  \item If predicted delta exceeds a safety threshold or degrades tyre life
    below the minimum margin: block and return an alternative.
\end{enumerate}

\subsection{Integration with Copilot}

The copilot calls the digital twin via the \mcp~gateway. The twin returns a
structured response with the predicted delta, confidence interval, and a
pass/fail decision. This response is included in the copilot recommendation
as a mandatory validation citation.
